\documentclass[10pt,conference,a4paper]{IEEEtran_EDM}
\IEEEoverridecommandlockouts

\usepackage{cite}
\usepackage{amsmath,amssymb,amsfonts}
\usepackage{algorithmic}
\usepackage{graphicx}
\usepackage{textcomp}
\usepackage{xcolor}
\usepackage{booktabs}
\usepackage{multirow}
\usepackage{url}
\usepackage[T1]{fontenc}

\def\BibTeX{{\rm B\kern-.05em{\sc i\kern-.025em b}\kern-.08em
    T\kern-.1667em\lower.7ex\hbox{E}\kern-.125emX}}
\def\confheader{}

\makeatletter
\newcommand{\linebreakand}{%
  \end{@IEEEauthorhalign}
  \hfill\mbox{}\par
  \mbox{}\hfill\begin{@IEEEauthorhalign}
}
\makeatother

\usepackage{flushend}
\begin{document}
\markboth{\confheader}{}
\title{Graph-Based vs Transition-Based Dependency Parsing for Uzbek}


\author{
\IEEEauthorblockN{Sanatbek Matlatipov}
\IEEEauthorblockA{\textit{Appl. mathematics and computer analysis}\\ 
\textit{National University of Uzbekistan}\\ 
Tashkent, Uzbekistan \\
s.matlatipov@nuu.uz}
\and
\IEEEauthorblockN{Maftuna Ismoilova}
\IEEEauthorblockA{\textit{ Innovative Development Agency}\\ 
\textit{MHESI}\\ 
Tashkent, Uzbekistan\\ 
m.ismoilova@ilmiy.uz}
\and
\IEEEauthorblockN{Mersaid Aripov}
\IEEEauthorblockA{\textit{Appl. mathematics and computer analysis}\\ 
\textit{National University of Uzbekistan}\\ 
Tashkent, Uzbekistan \\
mirsaid.aripov@nuu.uz}
}

\maketitle

\begin{abstract}
This article compares graph-based and transition-based neural dependency parsing architectures in Uzbek, an agglutinative Turkic language with few annotated resources. We examine two cutting-edge NLP frameworks: Stanza (graph-based DeepBiaffine parser) and spaCy (transition-based arc-eager parser), both of which are enhanced with TahrirchiBERT, a monolingual Uzbek BERT encoder. Experiments are carried out utilizing two existing Universal Dependencies treebanks containing 1,184 sentences from the fiction, academic, and news genres. Our results demonstrate a notable architecture-dependent trade-off: spaCy's joint end-to-end training produces superior POS tagging (UPOS 89.18 vs. 85.08), but Stanza's graph-based parser significantly outperforms in labeled dependency parsing (LAS 63.81 vs. 47.11).  Cross-treebank data augmentation improves POS accuracy consistently across both architectures but yields dramatically different LAS gains (+9.62 for Stanza vs.\ +1.76 for spaCy), suggesting that graph-based parsers exploit structural diversity in augmented data more effectively than transition-based alternatives. We further provide per-relation and per-feature error analyses that expose architecture-specific strengths, offering practical guidance for building NLP pipelines for agglutinative languages.
\end{abstract}

\begin{IEEEkeywords}
dependency parsing, Uzbek language, natural language processing, graph-based parsing, transition-based parsing, BERT, Universal Dependencies, morphosyntactic tagging
\end{IEEEkeywords}

\section{Introduction}

Uzbek is an agglutinative Turkic language spoken by over 35 million people, yet computational linguistic resources for Uzbek remain scarce compared to widely studied languages. The entire annotated dependency parsing data available for Uzbek consists of two Universal Dependencies (UD) treebanks~\cite{nivre2020universal} comprising only 1,184 sentences --- orders of magnitude smaller than resources available for English or even related Turkic languages such as Turkish~\cite{sulubacak2016universal}.

Dependency parsing --- the task of predicting the syntactic head and grammatical relation for each word in a sentence --- is fundamental for downstream NLP applications including information extraction, machine translation, and sentiment analysis. Two dominant neural architectures have emerged for this task: \textit{graph-based} parsers, which score all possible head--dependent pairs independently using biaffine attention~\cite{dozat2017deep}, and \textit{transition-based} parsers, which build dependency trees incrementally through a sequence of shift--reduce actions~\cite{nivre2008algorithms}.

While architectural comparisons have been conducted for related Turkic languages such as Turkish~\cite{eryigit2008dependency}, no study has evaluated modern neural frameworks with monolingual BERT encoders for Uzbek. Uzbek presents unique challenges: rich suffixal morphology encoding case, tense, person, and evidentiality; flexible SOV word order\cite{aripovkhakimov} with frequent center-embedding; and the availability of a monolingual BERT model (TahrirchiBERT) that may interact differently with each parsing architecture.

This paper makes three contributions:
\begin{enumerate}
    \item We present the \textbf{first cross-architecture comparison} of graph-based (Stanza~\cite{qi2020stanza}) and transition-based (spaCy~\cite{honnibal2020spacy}) dependency parsing for Uzbek, both augmented with the same TahrirchiBERT encoder.
    \item We quantify \textbf{architecture-dependent effects of cross-treebank data augmentation}, showing that merging two genre-complementary Uzbek UD treebanks produces dramatically different LAS improvements across parser types.
    \item We provide \textbf{detailed per-relation and per-feature error analyses} revealing that graph-based parsing excels at structurally complex relations while transition-based parsing achieves stronger local tagging.
\end{enumerate}

\section{Related Work}

\subsection{Dependency Parsing Architectures}

The graph-based DeepBiaffine parser~\cite{dozat2017deep} treats dependency parsing as a structured prediction problem, scoring all head--dependent pairs via biaffine attention over BiLSTM representations. Transition-based parsers~\cite{nivre2008algorithms} build trees incrementally using an arc-eager or arc-standard oracle. Comparative studies on CoNLL shared tasks have shown graph-based parsers generally achieve higher accuracy on languages with non-projective dependencies and rich morphology~\cite{de2017raw}, though transition-based parsers offer faster inference. Ballesteros et al.~\cite{ballesteros2015improved} showed that character composition models can improve transition-based parsing for morphologically rich languages by capturing subword structure.

\subsection{BERT-Augmented Parsing}

Pre-trained transformer models, especially BERT~\cite{devlin2019bert}, have substantially improved parsing accuracy. Kondratyuk and Straka~\cite{kondratyuk2019cross} demonstrated that multilingual BERT representations improve UD parsing across 75 languages. Stanza~\cite{qi2020stanza} originally used static word embeddings; our work extends it with monolingual BERT integration. spaCy~\cite{honnibal2020spacy} natively supports transformer backends through its \texttt{spacy-transformers} package.

\subsection{Turkic Language Parsing}

Eryiğit et al.~\cite{eryigit2008dependency} conducted foundational work on Turkish dependency parsing, demonstrating that morphology-aware sublexical representations significantly benefit parsing across both graph-based and transition-based strategies for agglutinative Turkic languages. For Uyghur, Zuo et al.~\cite{zuo2025mudt} proposed the MUDT treebank and reported 47.9\% annotation divergence between UD parsers and their scheme, highlighting cross-treebank challenges in Turkic languages. For Uzbek specifically, Elov~\cite{elov2025methods} evaluated biaffine graph parsers with UzBERT and XLM-R encoders on a larger 20k-sentence dataset derived from Uzbek national corpus texts, reporting LAS scores substantially higher than those achievable on our smaller UD-based datasets. However, that work did not compare parsing architectures (graph-based vs.\ transition-based) head-to-head, which is the focus of our study.

\subsection{Uzbek NLP Resources}

Recent work has also discussed methodological and practical issues in building a syntactic treebank for Uzbek outside the UD setting~\cite{uztreebank-esref}. Within the UD framework, two treebanks currently exist for Uzbek: UD\_Uzbek-UzUDT~\cite{matlatipov2025uzudt} containing 684 sentences from literary and academic texts, and UD\_Uzbek-UT~\cite{akhundjanova2025universal} containing 500 sentences from news and fiction. TahrirchiBERT is a monolingual Uzbek BERT-base model pre-trained on a diverse web corpus, providing 768-dimensional contextual embeddings.

\section{Datasets}

\subsection{UD Treebanks for Uzbek}

We use two publicly available UD treebanks for Uzbek:

\textbf{UD\_Uzbek-UzUDT}~\cite{matlatipov2025uzudt} contains 684 sentences (~7,800 tokens) from Uzbek literature and academic writing, manually annotated by six annotators (four linguists and two NLP engineers) on the INCEpTION platform with inter-annotator agreement exceeding 95\% for lemma and UPOS.

\textbf{UD\_Uzbek-UT}~\cite{akhundjanova2025universal} contains 500 sentences (~5,930 tokens) from news articles and fiction, created through semi-automatic annotation with full manual correction.

Neither treebank originally provided a standard three-way train/dev/test split. We re-split both using consistent UD proportions (~66\%/~7\%/~27\%), yielding the splits shown in Table~\ref{tab:splits}.

\begin{table}[htbp]
\caption{Data Splits for Uzbek UD Treebanks}
\begin{center}
\begin{tabular}{lccc}
\toprule
\textbf{Split} & \textbf{UzUDT} & \textbf{UT} & \textbf{Merged} \\
\midrule
Train & 451 & 330 & 781 \\
Dev & 45 & 33 & 78 \\
Test & 188 & 137 & 325 \\
\midrule
\textbf{Total} & \textbf{684} & \textbf{500} & \textbf{1,184} \\
\bottomrule
\end{tabular}
\label{tab:splits}
\end{center}
\end{table}

\subsection{Treebank Complementarity}

The two treebanks exhibit striking genre-driven differences that make them structurally complementary. Table~\ref{tab:complement} summarizes key distributional contrasts. UzUDT is rich in adverbial clauses (\texttt{advcl}: 444 vs.\ 106 in UT), reflecting literary converb constructions, while UT contains substantially more light-verb constructions (\texttt{compound:lvc}: 215 vs.\ 21) and named entity flat structures (\texttt{flat}: 106 vs.\ 11), typical of news text. The PROPN tag is 12$\times$ more frequent in UT (5.2\% vs.\ 0.3\%), reflecting the abundance of named entities in news.

\begin{table}[htbp]
\caption{Treebank Complementarity: Key Distributional Differences}
\begin{center}
\begin{tabular}{lrrl}
\toprule
\textbf{Feature} & \textbf{UzUDT} & \textbf{UT} & \textbf{Benefit} \\
\midrule
Genres & Fic, Acad & News, Fic & 3 genres \\
PROPN (\%) & 0.3 & 5.2 & NER coverage \\
\texttt{advcl} count & 444 & 106 & Clause struct. \\
\texttt{compound:lvc} & 21 & 215 & LVC patterns \\
Morph.\ features & 55 & 47 & Feature union \\
\bottomrule
\end{tabular}
\label{tab:complement}
\end{center}
\end{table}

Morphological annotation also differs: UzUDT annotates possessor agreement (\texttt{Number[psor]}, \texttt{Person[psor]}) and evidentiality (\texttt{Evident}=Fh/Nfh) --- typologically important Turkic features absent from UT. Conversely, UT includes abbreviation (\texttt{Abbr}), degree comparison (\texttt{Degree}), and finer numeral types (\texttt{Frac}, \texttt{Range}).

\section{Methodology}

We compare two fundamentally different neural parsing architectures, both augmented with TahrirchiBERT, a monolingual Uzbek BERT-base encoder.

\subsection{Stanza: Graph-Based Pipeline}

The Stanza pipeline~\cite{qi2020stanza} employs a two-stage architecture:

\textbf{Stage 1 --- POS Tagger:} A BiLSTM-based multi-task tagger jointly predicting UPOS, XPOS, and morphological features (UFeats). Input representations are either static FastText embeddings~\cite{bojanowski2017enriching} (300-dim) or TahrirchiBERT contextual embeddings (768-dim) with last-subword fusion.

\textbf{Stage 2 --- Dependency Parser:} A DeepBiaffine graph-based parser~\cite{dozat2017deep} scoring all head--dependent pairs via biaffine attention over BiLSTM representations, with separate MLP heads for arc prediction and relation labeling. The parser uses predicted POS tags from Stage~1 as input features.

\textbf{Subword Fusion:} TahrirchiBERT tokenizes text into WordPiece subwords whose boundaries rarely coincide with UD word tokens. For agglutinative Uzbek, where grammatical suffixes (case, tense, person) occupy the rightmost subwords, we employ \textit{last-subword selection} --- taking the hidden state of the final subword per UD token. For example, in \textit{bolalarning} (``of the children''), the subwords are \texttt{bola}, \texttt{\#\#lar}, \texttt{\#\#ning}, and the last subword \texttt{\#\#ning} carries the genitive case suffix targeted by UD annotation.

\subsection{spaCy: Transition-Based Pipeline}

The spaCy pipeline~\cite{honnibal2020spacy} uses a joint end-to-end architecture with four components trained simultaneously:

\textbf{Transformer:} TahrirchiBERT with strided spans (window=128, stride=96) providing contextual representations via \texttt{reduce\_mean} pooling.

\textbf{Tagger + Morphologizer:} TransformerListener components for XPOS prediction and joint UPOS + morphological feature prediction, respectively.

\textbf{Parser:} An arc-eager transition-based parser (\texttt{TransitionBasedParser.v2}) with hidden width 128 and maxout pieces 3.

Since spaCy does not natively support Uzbek, we created a custom Uzbek language class registered via entry points, enabling \texttt{spacy.blank("uz")} without patching spaCy source code. We also developed a CoNLL-U to spaCy DocBin converter preserving all UD annotation layers.

\subsection{Key Architectural Differences}

Table~\ref{tab:arch_diff} summarizes the fundamental differences between the two pipelines.

\begin{table}[htbp]
\caption{Architectural Comparison of Stanza and spaCy Pipelines}
\begin{center}
\begin{tabular}{lll}
\toprule
\textbf{Aspect} & \textbf{Stanza} & \textbf{spaCy} \\
\midrule
Parser type & Graph-based & Transition-based \\
Strategy & DeepBiaffine & Arc-eager \\
Training & Two-stage & Joint end-to-end \\
BERT fusion & Last-subword & Mean pooling \\
Max steps & 50,000 & 20,000 \\
\bottomrule
\end{tabular}
\label{tab:arch_diff}
\end{center}
\end{table}

\section{Experimental Setup}

\subsection{Experiment Matrix}

We train seven models across two frameworks, as summarized below:

\textbf{Stanza experiments:} E1.1/E1.2 (FastText, UzUDT/Merged), E2.1/E2.2 (TahrirchiBERT + last-subword, UzUDT/Merged), E3.1 (TahrirchiBERT + mean pooling, UzUDT).

\textbf{spaCy experiments:} S1.1/S1.2 (TahrirchiBERT, UzUDT/Merged).

For both frameworks, we use the recommended default hyperparameters from their respective documentation to ensure a fair ``out-of-the-box'' comparison reflecting the practical experience of NLP practitioners. Stanza models use Adam optimizer ($\beta$=(0.9, 0.999)), learning rate $3 \times 10^{-3}$, batch size 5,000 tokens, dropout 0.33, with early stopping on dev set. spaCy models use Adam with learning rate $5 \times 10^{-5}$, warmup over 250 steps, dropout 0.1, gradient accumulation 3, L2 weight decay 0.01, with patience of 1,600 steps. The differing hyperparameters (e.g., learning rates, dropout) reflect fundamentally different training paradigms: Stanza trains only the parser head atop frozen BERT features, while spaCy fine-tunes the full transformer jointly.

All experiments were conducted on a single NVIDIA RTX A6000 GPU (CUDA 12.4). Evaluation uses the standard UD CoNLL 2018 shared task metrics: UPOS, XPOS, UFeats (morphological features), UAS (unlabeled attachment score), and LAS (labeled attachment score).

\section{Results}

\subsection{Stanza Results}

Table~\ref{tab:stanza_results} presents the test-set results for all Stanza configurations.

\begin{table}[htbp]
\caption{Stanza Test Set Results (\%)}
\begin{center}
\setlength{\tabcolsep}{3.5pt}
\begin{tabular}{llccccc}
\toprule
\textbf{Run} & \textbf{Embed.} & \textbf{Data} & \textbf{UPOS} & \textbf{XPOS} & \textbf{UFeats} & \textbf{LAS} \\
\midrule
E1.1 & FastText & UzUDT & 79.19 & 79.81 & 66.61 & 51.24 \\
E1.2 & FastText & Merged & 80.26 & 83.20 & 66.98 & 62.40 \\
E2.1 & BERT$_\text{last}$ & UzUDT & 82.45 & 80.90 & 65.37 & 54.19 \\
\textbf{E2.2} & \textbf{BERT}$_\textbf{last}$ & \textbf{Merged} & \textbf{85.08} & \textbf{84.72} & \textbf{71.09} & \textbf{63.81} \\
E3.1 & BERT$_\text{mean}$ & UzUDT & 82.76 & 81.37 & 65.22 & 51.55 \\
\bottomrule
\end{tabular}
\label{tab:stanza_results}
\end{center}
\end{table}

The best Stanza configuration (E2.2) uses TahrirchiBERT with last-subword fusion on merged data, achieving UPOS 85.08\% and LAS 63.81\%. Comparing E2 (last-subword) vs.\ E3 (mean pooling) on UzUDT reveals that POS tagging performance is comparable, but last-subword fusion yields substantial parsing improvements (UAS +2.95, LAS +2.64), suggesting suffix-positional cues are more important for structural attachment decisions than for local tag assignment.

TahrirchiBERT consistently improves UPOS over FastText (+3.26 to +4.82 points). Notably, UFeats \textit{decreases} by 1.24 points with BERT on small data (E1.1~$\to$~E2.1) but \textit{increases} by 4.11 points on merged data (E1.2~$\to$~E2.2), indicating a minimum data threshold for contextual embeddings to learn fine-grained morphological feature bundles effectively.

Cross-treebank merging produces the single largest accuracy improvement: LAS jumps by +11.16 for FastText and +9.62 for BERT --- far exceeding the gains from upgrading embedding type.

\subsection{spaCy Results}

Table~\ref{tab:spacy_results} shows spaCy test-set performance.

\begin{table}[htbp]
\caption{spaCy Test Set Results (\%)}
\begin{center}
\begin{tabular}{lccccc}
\toprule
\textbf{Run} & \textbf{Data} & \textbf{UPOS} & \textbf{XPOS} & \textbf{UFeats} & \textbf{LAS} \\
\midrule
S1.1 & UzUDT & 86.50 & 86.72 & 50.55 & 45.35 \\
\textbf{S1.2} & \textbf{Merged} & \textbf{89.18} & \textbf{88.24} & \textbf{65.48} & \textbf{47.11} \\
\bottomrule
\end{tabular}
\label{tab:spacy_results}
\end{center}
\end{table}

\subsection{Cross-Architecture Comparison}

Table~\ref{tab:cross_arch} presents the central comparison between Stanza (E2.2) and spaCy (S1.2), both using TahrirchiBERT on merged data.

\begin{table}[htbp]
\caption{Cross-Architecture Comparison: Stanza vs spaCy (TahrirchiBERT, Merged Data, Test Set)}
\begin{center}
\begin{tabular}{lccr}
\toprule
\textbf{Metric} & \textbf{Stanza} & \textbf{spaCy} & \textbf{$\Delta$} \\
\midrule
UPOS & 85.08 & \textbf{89.18} & \textbf{+4.10} \\
XPOS & 84.72 & \textbf{88.24} & \textbf{+3.52} \\
UFeats & \textbf{71.09} & 65.48 & $-$5.61 \\
UAS & \textbf{72.39} & 66.81 & $-$5.58 \\
LAS & \textbf{63.81} & 47.11 & $-$\textbf{16.70} \\
\bottomrule
\end{tabular}
\label{tab:cross_arch}
\end{center}
\end{table}

A clear architecture-dependent trade-off emerges. spaCy achieves higher accuracy on POS tagging (UPOS +4.10, XPOS +3.52), likely because its joint end-to-end training objective reinforces tagging signals through parser loss gradients. However, Stanza's graph-based DeepBiaffine parser \textit{substantially} outperforms spaCy's arc-eager parser on labeled dependency parsing (LAS~63.81 vs.\ 47.11, a gap of 16.70 points). This is consistent with prior findings that exhaustive-search graph-based parsers excel on morphologically rich languages with frequent non-projective dependencies~\cite{de2017raw}.

\subsection{Data Augmentation Across Architectures}

Table~\ref{tab:augmentation} compares data augmentation effects (UzUDT~$\to$~Merged) across both architectures.

\begin{table}[htbp]
\caption{Effect of Cross-Treebank Data Augmentation by Architecture ($\Delta$ = Merged $-$ UzUDT)}
\begin{center}
\begin{tabular}{lrr}
\toprule
\textbf{Metric} & \textbf{Stanza $\Delta$} & \textbf{spaCy $\Delta$} \\
\midrule
UPOS & +2.63 & +2.68 \\
XPOS & +3.82 & +1.52 \\
UFeats & +5.72 & \textbf{+14.93} \\
UAS & +0.34 & $-$0.91 \\
LAS & \textbf{+9.62} & +1.76 \\
\bottomrule
\end{tabular}
\label{tab:augmentation}
\end{center}
\end{table}

Data augmentation is \textit{architecture-agnostic} for POS tagging: both frameworks gain approximately +2.6 UPOS points. However, the effects diverge dramatically for parsing and morphology. Stanza's LAS improves by +9.62 points while spaCy gains only +1.76, suggesting the graph-based parser exploits the structural diversity of merged data far more effectively. Conversely, spaCy's morphological accuracy jumps by +14.93 points (vs.\ +5.72 for Stanza), indicating that joint end-to-end training amplifies the benefit of additional labeled data for morphological features --- the morphologizer receives gradients from both tagger and parser objectives simultaneously.

\subsection{Per-Relation Error Analysis}

Fig.~\ref{fig:deprel_freq} shows the frequency distribution of the top 20 dependency relations in the combined corpus. Table~\ref{tab:per_relation} compares per-relation LAS F1 scores for selected relations from the spaCy S1.2 test evaluation, contextualizing where each architecture struggles.

\begin{figure}[htbp]
\centerline{\includegraphics[width=\columnwidth]{Figure2_DepRel_Frequency.png}}
\caption{Top 20 dependency relation frequencies in the combined Uzbek UD corpus. The distribution is heavily skewed toward \texttt{punct}, \texttt{obl}, and \texttt{nsubj}.}
\label{fig:deprel_freq}
\end{figure}

\begin{table}[htbp]
\caption{Per-Relation LAS F1 Scores (\%) — spaCy S1.2 (Merged Test Set)}
\begin{center}
\setlength{\tabcolsep}{4pt}
\begin{tabular}{lcccl}
\toprule
\textbf{Relation} & \textbf{P} & \textbf{R} & \textbf{F1} & \textbf{Note} \\
\midrule
\texttt{nummod} & 75.0 & 86.5 & 80.4 & Strongest \\
\texttt{root} & 72.9 & 74.5 & 73.7 & Core \\
\texttt{case} & 84.4 & 58.1 & 68.8 & P $\gg$ R \\
\texttt{obj} & 58.6 & 65.1 & 61.7 & \\
\texttt{nsubj} & 47.0 & 47.3 & 47.2 & \\
\texttt{obl} & 47.6 & 46.7 & 47.1 & \\
\texttt{compound:lvc} & 36.7 & 54.6 & 43.9 & LVCs \\
\texttt{nmod} & 27.0 & 34.2 & 30.2 & \\
\texttt{conj} & 25.8 & 14.9 & 18.9 & Weak \\
\bottomrule
\end{tabular}
\label{tab:per_relation}
\end{center}
\end{table}

The transition-based parser performs strongest on local, unambiguous relations (\texttt{nummod}: F1 80.4, \texttt{root}: 73.7) but struggles with structurally complex or long-distance relations (\texttt{conj}: 18.9, \texttt{nmod}: 30.2). The low \texttt{conj} F1 reflects coordination's inherent difficulty: conjoined elements can be arbitrarily distant and require identifying parallel syntactic structure across clause boundaries --- a challenge for greedy left-to-right processing. Similarly, \texttt{nmod} confusion arises from genitive case ambiguity in Uzbek, where possessive and attributive noun modification share the \textit{-ning} suffix. The graph-based parser, by scoring all head--dependent pairs exhaustively, better handles such structurally ambiguous relations.

\section{Discussion}

\subsection{Architecture Suitability}

Our results strongly favor graph-based parsing for Uzbek dependency analysis. The 16.70-point LAS gap between Stanza and spaCy is among the largest reported architecture differences in recent parsing literature and can be attributed to two primary factors: (1) Uzbek's SOV word order with frequent center-embedding creates many long-distance dependencies that arc-eager's greedy left-to-right processing handles poorly; (2) the graph-based parser's exhaustive search over all head--dependent pairs is better suited to the rich case system that disambiguates attachment sites.

Notably, Stanza achieves higher LAS \textit{despite} lower UPOS accuracy (85.08\% vs.\ 89.18\%), suggesting that raw POS tagging accuracy is not the primary driver of parsing quality. Rather, the parser architecture itself --- exhaustive biaffine scoring vs.\ incremental transition decisions --- is the dominant factor. spaCy's superior tagging arises from its joint end-to-end training, where parser loss gradients reinforce tagging signals, but this same joint objective may constrain the parser's capacity for complex structural decisions.

Our results complement Elov~\cite{elov2025methods}, who evaluated biaffine parsers with UzbBERT on a substantially larger dataset (~20k sentences). The different dataset sizes (781 vs.\ ~20k training sentences) and evaluation protocols preclude direct numerical comparison, but both studies confirm the effectiveness of graph-based biaffine parsing for Uzbek.

\subsection{Subword Fusion Strategy}

Comparing TahrirchiBERT with last-subword vs.\ mean pooling (E2.1 vs.\ E3.1), we find that the fusion strategy has negligible impact on POS tagging ($\Delta$UPOS = $-$0.31) but a substantial effect on parsing ($\Delta$LAS = +2.64 for last-subword). This suggests that suffix-level morphosyntactic cues preserved by last-subword selection are more critical for structural decisions (head selection, relation labeling) than for local POS classification, where distributional context from all subwords is equally informative.

\subsection{BERT--Data Interaction}

An important finding is the interaction between BERT and training data size. BERT's UFeats advantage flips from $-$1.24 on small data to +4.11 on merged data, a +5.35 amplification. This indicates that contextual embeddings require a minimum amount of training data to effectively learn fine-grained morphological feature bundles --- with too few examples, the model may overfit or fail to leverage the contextual signal.

\subsection{Limitations}

Several factors constrain absolute performance. First, the merged corpus (781 training sentences) is very small even by UD standards --- the median UD treebank contains approximately 10k sentences, placing Uzbek in the bottom 15\% by size. The test set of 325 sentences, while sufficient to detect differences of approximately 5 percentage points, limits the granularity of per-relation analyses. Second, BERT's WordPiece tokenization does not respect Uzbek morpheme boundaries; morpheme-aware tokenization could potentially improve results. Third, annotation inconsistencies between the two treebanks (UT has 92 UD validation warnings, differing morphological feature inventories) introduce label noise that particularly affects UFeats accuracy.

We acknowledge several methodological limitations: all results are from single training runs without confidence intervals or significance testing, which limits the strength of fine-grained conclusions given the small dataset. Additionally, our comparison is limited to two frameworks; evaluation of UDPipe~2.0, Trankit, or alternative BERT models (UzbBERT, XLM-R) would strengthen generalizability. These remain important directions for future work.

\section{Conclusion}

We presented the first cross-architecture comparison of graph-based and transition-based dependency parsing for Uzbek, revealing a pronounced trade-off: transition-based pipelines achieve superior POS tagging through joint end-to-end training, while graph-based parsers substantially outperform on labeled dependency parsing (+16.70 LAS). Cross-treebank data augmentation from genre-complementary sources is the single most impactful intervention, but its effectiveness is architecture-dependent --- graph-based parsers exploit structural diversity more effectively.

These findings provide practical guidance for Uzbek NLP pipeline design: graph-based architectures with monolingual BERT and linguistically-motivated subword fusion should be preferred when parsing accuracy is the primary objective. Future work includes evaluating alternative Uzbek BERT models (BERTbek), incorporating morpheme-aware tokenization, and expanding annotated resources to overcome the data bottleneck that remains the primary limiting factor.

All trained models are publicly available at \texttt{huggingface.co/Sanatbek/uzudt} and Github source code \texttt{SanatbekMatlatipov/robust-parsing-uzbek}.

\begin{thebibliography}{00}

\bibitem{nivre2020universal} J. Nivre, M.-C. de Marneffe, F. Ginter, J. Hajič, C. D. Manning, S. Pyysalo, S. Schuster, F. Tyers, and D. Zeman, "Universal Dependencies v2: An evergrowing multilingual treebank collection," in \textit{Proceedings of the Twelfth Language Resources and Evaluation Conference}, Marseille, France, 2020, pp. 4034--4043, ISBN: 979-10-95546-34-4. [Online]. Available: \url{https://aclanthology.org/2020.lrec-1.497/}

\bibitem{sulubacak2016universal}
U. Sulubacak, M. Gokirmak, F. Tyers, Ç. Çöltekin, J. Nivre, and G. Eryiğit, "Universal Dependencies for Turkish," in \textit{Proceedings of COLING 2016, the 26th International Conference on Computational Linguistics: Technical Papers}, Osaka, Japan, Dec. 2016, pp. 3444--3454. [Online]. Available: \url{https://aclanthology.org/C16-1325/}


\bibitem{dozat2017deep}
T. Dozat and C. D. Manning, "Deep biaffine attention for neural dependency parsing," \textit{arXiv preprint} arXiv:1611.01734, 2016, doi: 10.48550/arXiv.1611.01734. [Online]. Available: \url{https://arxiv.org/abs/1611.01734}

\bibitem{nivre2008algorithms}
J. Nivre, "Algorithms for deterministic incremental dependency parsing," \textit{Computational Linguistics}, vol. 34, no. 4, pp. 513--553, 2008, doi: 10.1162/coli.07-056-R1-07-027. [Online]. Available: \url{https://aclanthology.org/J08-4003/}

\bibitem{eryigit2008dependency} Eryiğit, G., Nivre, J., Oflazer, K. (2018). Dependency Parsing of Turkish. In: Oflazer, K., Saraçlar, M. (eds) Turkish Natural Language Processing. Theory and Applications of Natural Language Processing. Springer, Cham. https://doi.org/10.1007/978-3-319-90165-7\_7

\bibitem{aripovkhakimov} Aripov, M., Khakimov, M., Matlatipov, S., Sirojiddinov, Z. (2022). Analysis and Processing of the Uzbek Language on the Multi-language Modelled Computer Translator Technology. In: Vetulani, Z., Paroubek, P., Kubis, M. (eds) Human Language Technology. Challenges for Computer Science and Linguistics. LTC 2019. Lecture Notes in Computer Science(), vol 13212. Springer, Cham. https://doi.org/10.1007/978-3-031-05328-3\_6

\bibitem{qi2020stanza}
P. Qi, Y. Zhang, Y. Zhang, J. Bolton, and C. D. Manning, "Stanza: A Python natural language processing toolkit for many human languages," in \textit{Proceedings of the 58th Annual Meeting of the Association for Computational Linguistics: System Demonstrations}, Online, Jul. 2020, pp. 101--108, doi: 10.18653/v1/2020.acl-demos.14. [Online]. Available: \url{https://aclanthology.org/2020.acl-demos.14/}

\bibitem{honnibal2020spacy} Honnibal, M., Montani, I., Van Landeghem, S., \& Boyd, A. (2020). spaCy: Industrial-strength Natural Language Processing in Python. https://doi.org/10.5281/zenodo.1212303.

\bibitem{de2017raw}
M. de Lhoneux, Y. Shao, A. Basirat, E. Kiperwasser, S. Stymne, Y. Goldberg, and J. Nivre, "From raw text to Universal Dependencies---look, no tags!," in \textit{Proceedings of the CoNLL 2017 Shared Task: Multilingual Parsing from Raw Text to Universal Dependencies}, Vancouver, Canada, Aug. 2017, pp. 207--217, doi: 10.18653/v1/K17-3022. [Online]. Available: \url{https://aclanthology.org/K17-3022/}

\bibitem{ballesteros2015improved}
M. Ballesteros, C. Dyer, and N. A. Smith, "Improved transition-based parsing by modeling characters instead of words with LSTMs," in \textit{Proceedings of the 2015 Conference on Empirical Methods in Natural Language Processing}, Lisbon, Portugal, Sep. 2015, pp. 349--359, doi: 10.18653/v1/D15-1041. [Online]. Available: \url{https://aclanthology.org/D15-1041/}

\bibitem{devlin2019bert}
J. Devlin, M.-W. Chang, K. Lee, and K. Toutanova, "BERT: Pre-training of deep bidirectional transformers for language understanding," in \textit{Proceedings of the 2019 Conference of the North American Chapter of the Association for Computational Linguistics: Human Language Technologies, Volume 1 (Long and Short Papers)}, Minneapolis, MN, USA, Jun. 2019, pp. 4171--4186, doi: 10.18653/v1/N19-1423. [Online]. Available: \url{https://aclanthology.org/N19-1423/}

\bibitem{kondratyuk2019cross}
D. Kondratyuk and M. Straka, "75 languages, 1 model: Parsing Universal Dependencies universally," in \textit{Proceedings of the 2019 Conference on Empirical Methods in Natural Language Processing and the 9th International Joint Conference on Natural Language Processing (EMNLP-IJCNLP)}, Hong Kong, China, Nov. 2019, pp. 2779--2795, doi: 10.18653/v1/D19-1279. [Online]. Available: \url{https://aclanthology.org/D19-1279/}

\bibitem{zuo2025mudt} J. Zuo, Y. Wang, and Y. Pan, ``Modern Uyghur Dependency Treebank (MUDT): An integrated morphosyntactic framework for a low-resource language,'' \textit{arXiv:2507.21536}, 2025.

\bibitem{elov2025methods} B. B. Elov, ``Methods and algorithms for dependency parse of texts of the national Uzbekistan corpus,'' in \textit{Proc. Sci. Conf. Uzbek Lang. Literature}, Tashkent, Uzbekistan, 2025, doi: 10.52773/tsuull.conf.2025/wfwm9571.

\bibitem{uztreebank-esref}
E. Adalı and A. O. Xolmo'minovna, "UzTreebank: Methodological and practical issues in building a syntactic treebank for the Uzbek language," in \textit{2025 10th International Conference on Computer Science and Engineering (UBMK)}, Istanbul, Turkiye, 2025, pp. 445--450, doi: 10.1109/UBMK67458.2025.11206779.

\bibitem{matlatipov2025uzudt}
S. Matlatipov and E. Kuriyozov, ``UD Uzbek UzUDT,'' \textit{Universal Dependencies}, release 2.17, 2025. [Online]. Available: \url{https://universaldependencies.org/treebanks/uz_uzudt/index.html}

\bibitem{akhundjanova2025universal}
A. Akhundjanova and L. Talamo, "Universal Dependencies treebank for Uzbek," in \textit{Proceedings of the Third Workshop on Resources and Representations for Under-Resourced Languages and Domains (RESOURCEFUL-2025)}, Tallinn, Estonia, Mar. 2025, pp. 1--6, ISBN: 978-9908-53-121-2. [Online]. Available: \url{https://aclanthology.org/2025.resourceful-1.1/}

\bibitem{bojanowski2017enriching}
P. Bojanowski, E. Grave, A. Joulin, and T. Mikolov, "Enriching word vectors with subword information," \textit{Transactions of the Association for Computational Linguistics}, vol. 5, pp. 135--146, 2017, doi: 10.1162/tacl\_a\_00051. [Online]. Available: \url{https://aclanthology.org/Q17-1010/}

\end{thebibliography}

\end{document}
